% Part: applied-modal-logics
% Chapter: temporal-logic
% Section: language-epistemic-logic

\documentclass[../../../include/open-logic-section]{subfiles}

\begin{document}

\olfileid[pt]{aml}{tl}{sem}

\olsection{Semântica da lógica temporal}

\begin{defn}
A linguagem básica da lógica temporal contém:
\begin{enumerate}
  \tagitem{prvFalse}{A constante proposicional de !!{falsity}~$\lfalse$.}{}
  \tagitem{prvTrue}{A constante proposicional de !!{truth}~$\ltrue$.}{}
  \item Um conjunto de elementos !!{denumerable}s, as !!{propositional variable}s: $\Obj
    p_0$, $\Obj p_1$, $\Obj p_2$, \dots
  \item Os conectivos proposicionais: \startycommalist
  \iftag{prvNot}{\ycomma $\lnot$ (negação)}{}%
  \iftag{prvAnd}{\ycomma $\land$ (conjunção)}{}%
  \iftag{prvOr}{\ycomma $\lor$ (disjunção)}{}%
  \iftag{prvIf}{\ycomma $\lif$ (!!{conditional})}{}%
  \iftag{prvIff}{\ycomma $\liff$ (!!{biconditional})}{}.
  \item Os operadores de passado $\Ptemp$ e $\Htemp$.
  \item Os operadores de futuro $\Ftemp$ e $\Gtemp$.
\end{enumerate}
\end{defn}

Mais adiante, discutiremos o possível acréscimo de outros tipos de operadores modais.

\begin{defn}
As \emph{!!^{formula}s} da linguagem temporal são definidas
  indutivamente da seguinte maneira:
\begin{enumerate}
\tagitem{prvFalse}{$\lfalse$ é uma !!{formula} atômica.}{}

\tagitem{prvTrue}{$\ltrue$ é uma !!{formula} atômica.}{}

\item Toda variável proposicional $\Obj p_i$ é uma !!{formula} (atômica).

\tagitem{prvNot}{Se $!A$ é !!a{formula}, então $\lnot !A$ é
  !!a{formula}.}{}

\tagitem{prvAnd}{Se $!A$ e $!B$ são !!{formula}s, então $(!A \land
  !B)$ é !!a{formula}.}{}

\tagitem{prvOr}{Se $!A$ e $!B$ são !!{formula}s, então $(!A \lor !B)$
  é !!a{formula}.}{}

\tagitem{prvIf}{Se $!A$ e $!B$ são !!{formula}s, então $(!A \lif !B)$
  é !!a{formula}.}{}

\tagitem{prvIff}{Se $!A$ e $!B$ são !!{formula}s, então $(!A \liff !B)$
  é !!a{formula}.}{}

\item Se $!A$ é !!a{formula}, então $\Ptemp  !A$, $\Htemp  !A$, $F !A$, $\Gtemp  !A$ são todas
  !!{formula}s.

\tagitem{limitClause}{Nada mais é !!a{formula}.}{}
\end{enumerate}
\end{defn}

A semântica das lógicas temporais é dada em termos de modelos relacionais, como ocorre com outros tipos de lógicas intensionais. 

\begin{defn}
  Um \emph{modelo} para a linguagem temporal é uma tripla
  $\mModel{M} = \tuple{T, \prec, V}$, em que
  \begin{enumerate}
  \item $T$ é um conjunto não vazio, interpretado como instantes temporais.
  \item $\prec$ é uma relação binária sobre $T$.
  \item $V$ é uma função que atribui a cada !!{propositional
    variable}~$p$ um conjunto $V(p)$ de instantes temporais.
  \end{enumerate}
  Quando $t \prec t'$, dizemos que $t$ \emph{precede}~$t'$. 
  Quando $t \in V(p)$, dizemos que $p$ é \emph{verdadeira em}~$t$.
\end{defn}

Por enquanto, não impomos nenhuma condição à relação de precedência $\prec$.
Isso significa que, neste momento, não há restrições à estrutura de nossos
modelos temporais: podemos ter modelos nos quais o tempo seja linear,
ramificado, circular ou possua qualquer outra estrutura. 

Assim como na lógica modal normal, todo modelo temporal determina quais
!!{formula}s são verdadeiras em quais de seus instantes. Usamos a mesma
notação ``o modelo $\mModel{M}$ torna a !!{formula}~$!A$ verdadeira no
instante~$t$'' para a noção básica da semântica relacional. A relação é
definida indutivamente e, para todos os operadores não modais, é idêntica
à do caso modal normal.

\begin{defn}\ollabel{defn:tmodels}
  A \emph{verdade de !!a{formula}~$!A$ em~$t$} em um~$\mModel M$, em símbolos
  $\mSat{M}{!A}[t]$, é definida indutivamente da seguinte maneira:
  \begin{enumerate}
  \tagitem{prvFalse}{\indcase{!A}{\lfalse}{Nunca $\mSat{M}{\lfalse}[t]$}.}{}
  \tagitem{prvTrue}{\indcase{!A}{\ltrue}{Sempre $\mSat{M}{\ltrue}[t]$}.}{}
  \item $\mSat{M}{p}[t]$ sse $t \in V(p)$
  \tagitem{prvNot}{\indcase{!A}{\lnot !B}{$\mSat{M}{\indfrm}[t]$ sse
    $\mSat/{M}{!B}[t]$}.}{}
  \tagitem{prvAnd}{\indcase{!A}{(!B \land !C)}{$\mSat{M}{\indfrm}[t]$ sse
    $\mSat{M}{!B}[t]$ e $\mSat{M}{!C}[t]$}.}{}
  \tagitem{prvOr}{\indcase{!A}{(!B \lor !C)}{$\mSat{M}{\indfrm}[t]$ sse
    $\mSat{M}{!B}[t]$ ou $\mSat{M}{!C}[t]$} (ou ambos).}{}
  \tagitem{prvIf}{\indcase{!A}{(!B \lif !C)}{$\mSat{M}{\indfrm}[t]$ sse
    $\mSat/{M}{!B}[t]$ ou $\mSat{M}{!C}[t]$}.}{}
  \tagitem{prvIff}{\indcase{!A}{(!B \liff !C)}{$\mSat{M}{\indfrm}[t]$ sse
    ou ambos $\mSat{M}{!B}[t]$ e $\mSat{M}{!C}[t]$, ou
    nem $\mSat{M}{!B}[t]$ nem $\mSat{M}{!C}[t]$}.}{}
  \item\ollabel{defn:sub:mmodels-p}
    \indcase{!A}{\Ptemp  !B}{$\mSat{M}{\indfrm}[t]$ sse
    $\mSat{M}{!B}[t']$ para algum $t' \in T$ tal que $t' \prec t$}
\item\ollabel{defn:sub:mmodels-h}
    \indcase{!A}{\Htemp  !B}{$\mSat{M}{\indfrm}[t]$ sse
    $\mSat{M}{!B}[t']$ para todo $t' \in T$ tal que $t' \prec t$}
   \item\ollabel{defn:sub:mmodels-f}
    \indcase{!A}{\Ftemp  !B}{$\mSat{M}{\indfrm}[t]$ sse
    $\mSat{M}{!B}[t']$ para algum $t' \in T$ tal que $t \prec t'$}
\item\ollabel{defn:sub:mmodels-g}
    \indcase{!A}{\Gtemp  !B}{$\mSat{M}{\indfrm}[t]$ sse
    $\mSat{M}{!B}[t']$ para todo $t' \in T$ tal que $t \prec t'$}
  \end{enumerate} 
\end{defn}

Com base na semântica, pode-se observar que os operadores $\Ptemp$ e~$\Htemp$ são duais, assim como os
operadores $\Ftemp$ e $\Gtemp$; portanto, poderíamos definir $\Htemp  !A$ como $\lnot \Ptemp \lnot !A$, e fazer o mesmo com $\Gtemp$ e~$\Ftemp$.

\end{document}
